\section{总结与延伸}

这一讲的目标，其实是把数据从”看不见的前置条件”变成”显式的工程对象”。数据不是被动存在的输入，而是可以通过来源选择、过滤、混合、再写、再对齐被系统性塑造的。

\begin{importantbox}{最后的结论}
\begin{enumerate}
    \item 数据工作本身就是一个长期、并行、可扩展、且高度依赖经验的工程。
    \item 预训练、mid-training、post-training 的数据目标不同，不能混为一谈。
    \item 网页语料不是天然可用，必须经过抓取、清洗、去重、质量打分和合规检查。
    \item 任务化数据、指令数据和合成数据已经成为后训练的重要组成。
    \item 版权和平台条款不是边角问题，而是训练数据供应链的一部分。
    \item 数据管线中的大量步骤仍然是启发式的，有大量改进空间。
\end{enumerate}
\end{importantbox}

讲者想传达的最朴素一句话是：\textbf{模型能力不是”长出来”的，而是”喂出来”的。} 你怎么喂，决定你得到什么样的模型。

\subsection{本讲的核心问题链}

回顾这一讲的逻辑链条：
\begin{enumerate}
    \item \textbf{我们用什么数据训练？} 从 Wikipedia + 书籍到数十万亿 token 的网页数据，来源越来越广、清洗越来越精。
    \item \textbf{我们如何过滤、混合和验证数据？} 从简单规则到分类器、从精确去重到模糊去重、从人工策划到基准化评测。
    \item \textbf{有什么法律和伦理约束？} 版权、许可、ToS、隐私——每一层都可能独立构成限制。
    \item \textbf{这些选择如何变成你得到的模型？} 数据配方决定模型能力分布，数据质量决定模型能力上限。
\end{enumerate}

\subsection{拓展阅读}

\begin{itemize}
    \item Stanford CS324 课程中的 data 和 legality 章节提供了更多细节
    \item Longpre et al., \textit{A Pretrainer's Guide to Training Data} (2023) 系统梳理了预训练数据的各个方面
    \item Penedo et al., \textit{The FineWeb Datasets} (2024) 是现代网页数据清洗的最佳实践参考
    \item Li et al., \textit{DataComp-LM} (2024) 定义了数据处理方法的标准化评测框架
\end{itemize}

\end{document}
